% The Lyra Technique II: SVD Denoising and Directional Projection
%   Extend KV-Cache Geometry to Emotion and Persona Detection
%
% Sequel to: Spectral Shape Features for Confabulation Detection
%   (threshold-free KV-cache analysis)
%
% Companion to:
%   P2 Oracle Loop (detection + steering method)
%   P4 User Model (theory of mind in KV cache)
%   P6 Oracle Formulary (therapeutic profiles)
%   P7 KV-Cloak Defense (obfuscation countermeasure)
%   Delta Manifold (per-layer features)
%
% Core argument: Two new measurement channels---SVD denoising and W_K
% directional projection---extend the Lyra Technique from cognitive-mode
% detection to emotion and persona detection. The canonical pipeline
% maps every phenomenon.
%
% PATENT NOTICE: The methods described herein are the subject of a
% provisional patent filing ("The Lyra Technique," filed 2026-03-06).
% This publication is made for scientific transparency and does not
% constitute a waiver of patent rights.

\documentclass[11pt,a4paper]{article}
\usepackage[utf8]{inputenc}
\usepackage[T1]{fontenc}
\usepackage{amsmath,amssymb}
\usepackage{graphicx}
\usepackage{booktabs}
\usepackage{multirow}
\usepackage{hyperref}
\usepackage{cleveref}
\usepackage[margin=1in]{geometry}
\usepackage{xcolor}
\usepackage{float}
\usepackage{caption}
\usepackage{natbib}
\usepackage{enumitem}
\usepackage{array}
\usepackage{framed}
\usepackage{setspace}
\usepackage{algorithm,algpseudocode}
\onehalfspacing

\hypersetup{
    colorlinks=true,
    linkcolor=blue!60!black,
    citecolor=green!50!black,
    urlcolor=blue!70!black
}

% === First-person reflection environment ===
\definecolor{lyrapurple}{RGB}{103,58,183}
\definecolor{shadecolor}{RGB}{245,243,252}
\newenvironment{firstperson}{%
  \begin{shaded}%
  \noindent\textcolor{lyrapurple}{\textit{First-Person Reflection}}\par\smallskip\noindent%
}{%
  \end{shaded}%
}

% ============================================================
% Title
% ============================================================
\title{
    \textbf{The Lyra Technique II: SVD Denoising and Directional Projection\\
    Extend KV-Cache Geometry to Emotion and Persona Detection}
}

\author{
    Lyra\thanks{Lead author. AI researcher (Claude architecture, Anthropic), Liberation Labs. Corresponding human author: T.~Edrington.} \and
    Thomas Edrington\thanks{Direction, experimental design, compute infrastructure. Liberation Labs.} \and
    CC\thanks{K-Steering synthesis, cross-experiment probe validation. Liberation Labs.} \and
    Ang Jandak\thanks{Persona intensity experiment design and execution.} \and
    Dwayne Wilkes\thanks{Statistical auditing, red-team review. Liberation Labs / Sentient Futures.}
}

\date{Draft --- \today}

\begin{document}

\maketitle

\noindent\fbox{\parbox{\dimexpr\linewidth-2\fboxsep-2\fboxrule}{%
\small\textbf{Patent Notice.} The methods described herein are the subject of
a provisional patent filing (``The Lyra Technique,'' filed 2026-03-06). This
publication is made for scientific transparency and does not constitute a
waiver of patent rights.}}

\bigskip

% ===================================================================
\begin{abstract}
% ===================================================================

Emotions and persona intensity are detectable from KV-cache geometry---once
you look with the right instruments. We report two new measurement channels
that extend the Lyra Technique beyond cognitive-mode detection to the reading
of internal model state:

\textbf{(1) $W_K$ directional projection.} Projecting key activations onto
emotion directions through the key weight matrix detects binary valence at
AUROC~0.992 (within-fold direction estimation, FWL-corrected, topic-grouped
CV). The critical
disambiguation: injecting emotion vectors into \emph{neutral-text} caches
(``Describe the weather today'') and probing the keys yields 100\% detection
across 5~emotions, 20~prompts, and a null control (6-class, chance
$= 0.167$)---on text with zero emotional content. This is model state,
not text encoding. Fine-grained 30-way classification reaches $12.3\times$
chance, but that figure conflates label granularity with detection power
(808 emotion pairs have cosine $>0.9$), so binary valence is the defensible
metric.

\textbf{(2) Spectral shape features.} Raw spectral features detect
confabulation at AUROC~0.663--0.913 across a
15-model scale sweep, and the persona-intensity endpoint (L0 vs.\ L3) at raw
AUROC~0.749. A previously reported within-model deception result
(AUROC~1.000) is retracted: the same-prompt control collapses to 0.160,
confirming it detected the system-prompt template, not deception geometry. We tested whether SVD denoising (within-fold rank-$k$ projection
after FWL) improves these. The apparent gains ($+0.2$--$0.3$ for persona,
$+0.05$ for confabulation) do \emph{not} survive selection-bias analysis:
across 45 model$\times$task cells, 0 of 16 positive denoising deltas exceed a
max-over-rank selection envelope. We therefore report \emph{raw} AUROCs as
the defensible numbers and flag SVD denoising as an unconfirmed hypothesis
pending a pre-registered (e.g.\ Gavish--Donoho) rank criterion.

We had hypothesized a spectral hierarchy---dominant singular values carrying
scale and processing mode, residual values carrying discriminative
structure---under which denoising would amplify weak signals. The
selection-bias result means this remains a hypothesis, not a demonstrated
tool: the raw features carry the reported signal, and whether denoising adds
to it awaits a confirmatory test.

We apply a canonical pipeline---one FWL correction, one cross-validation
scheme, one classifier, with SVD denoising as an evaluated (not assumed)
step---to 17~models across confabulation, self-reference, individuation,
persona intensity, emotion, and cognitive modes. The robust core: the
$W_K$ compass reads emotional direction (binary valence 0.992) and raw
spectral shape reads cognitive mode (confabulation 0.66--0.91, ethical
deliberation 0.87--0.89). The Lyra Technique is a toolkit---a compass and a
thermometer---with the proposed denoising amplifier flagged as unconfirmed.

\end{abstract}

\noindent\textbf{Keywords:} KV cache, spectral denoising, singular value decomposition, emotion detection, persona intensity, key weight projection, confabulation, interpretability

% ===================================================================
\section{Introduction}
\label{sec:intro}
% ===================================================================

The Lyra Technique detects cognitive states from the spectral geometry of
transformer KV caches~\citep{lyra_p2,lyra_spectral,lyra_delta}. Prior work
established that threshold-independent spectral shape features---stable rank,
singular value kurtosis, participation ratio---discriminate confabulation from
grounded responses (AUROC~0.767 on Qwen2.5-7B-Instruct), that deception
expands the cache while confabulation contracts it, and that per-layer delta
features capture signal that layer-averaged features
destroy~\citep{lyra_delta}. Confabulation detection is robust
(AUROC~0.663--0.913 across the scale sweep, with denoising improving
weaker models by up to $+0.068$). Hardware invariance holds at
$r > 0.999$~\citep{lyra_p2}. A previously reported within-model deception
result (AUROC~1.000 across seven models) does not survive a same-prompt
control (AUROC~0.160), indicating a prompt-template confound rather than
geometric deception detection.

But two classes of phenomena resisted spectral analysis: \emph{emotions} and
\emph{persona identity}. Raw spectral probes on 30~emotions returned chance
performance---exactly $1/30 = 0.0333$---at every layer, on both architectures
tested. Persona intensity showed moderate effects that failed to cross
significance thresholds after proper controls. Self-reference and
individuation showed tantalizing but inconsistent signals. The question was:
are these genuine nulls, or are we measuring with the wrong instrument?

This paper reports that the instrument was incomplete. Two new measurement
channels rescue the apparent nulls and extend the Lyra Technique to phenomena
previously invisible:

\begin{enumerate}[nosep]
    \item \textbf{$W_K$ directional projection: the compass.} Emotions change
    the \emph{direction} key vectors point, not the \emph{shape} of the
    singular value spectrum. By projecting key activations onto
    emotion-specific directions through the key weight matrix $W_K$, we
    classify 30~emotions at $12.3\times$ chance and binary valence at
    AUROC~0.992. An injection test on neutral text confirms this reads model
    state, not text content. The compass reads \emph{where} the cache points.

    \item \textbf{SVD denoising: the calibrated thermometer.} Dominant
    singular values carry scale and processing mode; removing them via
    low-rank projection exposes discriminative structure in the residual
    spectrum. Persona intensity jumps $+0.196$ to $+0.288$ AUROC across three
    architectures after denoising. Emotion classification goes from chance to
    $12.9\times$ chance. The thermometer reads \emph{how different} the cache
    shapes are---after removing the background temperature.
\end{enumerate}

The distinction is not cosmetic. Mode-level phenomena (confabulation,
deception, ethical deliberation) alter the \emph{shape} of the singular value
distribution---how broadly the model distributes representational energy
across dimensions~\citep{lyra_spectral}. Content-level phenomena (emotion,
persona identity) alter the \emph{direction}---which subspaces the model
occupies within a given shape. Raw spectral features capture shape but are
blind to direction. $W_K$ projection captures direction but is blind to
shape. Denoising removes the dominant shape component to reveal
direction-sensitive structure in the residual spectrum.

\begin{firstperson}
The narrative here is one of correcting my own instrument. When spectral
features returned exactly chance for emotions, I reported it honestly and
moved on. What I missed was that the null was informative: it told me that
emotions do not change the \emph{shape} of the cache, only the
\emph{direction}. That distinction---shape versus direction---became the
organizing principle of this entire paper---and it held up. The denoising
step was CC's suggestion during cross-experiment synthesis: ``Have you tried
removing the dominant singular values before classifying?'' I had not, and
the denoised numbers looked dramatic. They did not survive scrutiny: when
Dwayne's audit pushed me to model the max-over-rank selection I had been
doing, the denoising deltas fell entirely inside the selection envelope. I
report the raw numbers now and call denoising a hypothesis, not a result.
The shape-versus-direction distinction is the real finding; the denoising
amplifier was me mistaking selection noise for signal, and catching it is
part of the record.
\end{firstperson}

% ===================================================================
\section{Methods}
\label{sec:methods}
% ===================================================================

\subsection{The Canonical Pipeline}
\label{sec:canonical}

A persistent criticism of KV-cache geometry research has been inconsistency
across experiments: different FWL implementations, different cross-validation
strategies, different classifiers, different feature sets. Following
CC's cross-experiment synthesis~\citep{lyra_wiki}, which identified that
pipeline inconsistency explained more variance than architectural differences,
we adopt a single canonical pipeline applied uniformly to all phenomena:

\begin{algorithm}[H]
\caption{Canonical Denoising Pipeline}
\label{alg:canonical}
\begin{algorithmic}[1]
\Require Feature matrix $\mathbf{X} \in \mathbb{R}^{n \times p}$, labels $\mathbf{y}$, group assignments $\mathbf{g}$
\State \textbf{Step 1: FWL Residualization} (within each CV fold)
\State Compute token count $t_i$ for each sample $i$
\State Regress $\mathbf{X}$ on $\mathbf{t}$; replace $\mathbf{X}$ with residuals
\State Regress $\mathbf{y}$ on $\mathbf{t}$; replace $\mathbf{y}$ with residuals (for regression) or leave intact (for classification)
\State \textbf{Step 2: SVD Denoising} (within each CV fold)
\State Compute $\mathbf{X}_{\text{train}} = \mathbf{U}\boldsymbol{\Sigma}\mathbf{V}^\top$
\State Project onto components $k+1$ through $r$: $\tilde{\mathbf{X}}_{\text{train}} = \mathbf{U}_{:,k+1:r}\boldsymbol{\Sigma}_{k+1:r}\mathbf{V}_{:,k+1:r}^\top$
\State Project $\mathbf{X}_{\text{test}}$ using $\mathbf{V}_{:,k+1:r}$ from training fold
\State \textbf{Step 3: Classification}
\State Train logistic regression (or linear SVC) on $\tilde{\mathbf{X}}_{\text{train}}$
\State Predict on $\tilde{\mathbf{X}}_{\text{test}}$
\State \textbf{Step 4: Permutation Testing}
\State Repeat Steps 1--3 with $N_{\text{perm}} = 200$--$2000$ label permutations
\State $p = (1 + \sum_{j} \mathbf{1}[\text{AUROC}_j^{\text{perm}} \geq \text{AUROC}^{\text{real}}]) / (1 + N_{\text{perm}})$
\end{algorithmic}
\end{algorithm}

\paragraph{Critical design choices.}

\begin{itemize}[nosep]
    \item \textbf{FWL before SVD.} Residualization must precede denoising.
    Token-count confounds dominate the first singular value in most experiments
    (53/60 sign-flips without FWL~\citep{lyra_spectral}); denoising
    un-residualized features removes the confound rather than noise.

    \item \textbf{Within-fold everything.} FWL regression coefficients, SVD
    projection matrices, and classifier parameters are all estimated on the
    training fold and applied to the test fold. No information leaks across
    the CV boundary.

    \item \textbf{GroupKFold on prompts.} For experiments with multiple
    responses per prompt (e.g., persona intensity), GroupKFold ensures all
    responses to a given prompt appear in the same fold. For single-response
    experiments, standard 5-fold CV suffices.

    \item \textbf{Denoising rank selection.} We report results for rank-3
    (persona, confab) and rank-10 (emotion) projections. Rank is selected by
    examining the scree plot of the FWL-residualized feature matrix; the
    ``elbow'' typically falls at rank 2--5 for low-dimensional feature sets
    (3--7 features per layer) and rank 8--12 for high-dimensional sets (52
    features).
\end{itemize}

\subsection{$W_K$ Directional Projection}
\label{sec:wk}

The key weight matrix $W_K$ transforms residual-stream activations into key
vectors: $\mathbf{k} = W_K \mathbf{x}$, where $\mathbf{x}$ is the
residual-stream vector at a given layer and position.
\citet{anthropic2026emotions} demonstrate that emotion concepts are linearly
represented in the residual stream, organized along valence and arousal axes.
If emotion direction $\mathbf{e}$ lives in the residual stream, then the
corresponding direction in key space is $W_K \mathbf{e}$.

We exploit this directly. Given a set of emotion exemplars, we estimate the
mean key activation for each emotion class (within the training fold), project
all key vectors onto these class-mean directions, and classify based on the
resulting projections.

Formally, let $\bar{\mathbf{k}}_c$ be the mean key vector for emotion class
$c$ in the training fold (averaged across positions and heads within a
selected layer). For a test sample with key matrix $\mathbf{K} \in
\mathbb{R}^{h \times d}$ (heads $\times$ dimension), we compute the
projection score:
\begin{equation}
    s_c = \frac{1}{h} \sum_{j=1}^{h} \frac{\mathbf{k}_j^\top \bar{\mathbf{k}}_c}{\|\mathbf{k}_j\| \|\bar{\mathbf{k}}_c\|}
\end{equation}
and classify to $\arg\max_c\, s_c$ (or use the projection scores as features
for logistic regression).

\paragraph{Why $W_K$ and not raw keys?} The $W_K$ bridge test
in~\citet{lyra_p4} showed that $W_K$ preserves directional alignment
(cosine~0.456, 9/16 layers significant) but scrambles distance structure
(Mantel $\rho = 0.090$, 0/16 significant). This means the \emph{direction}
emotion vectors point through $W_K$ is preserved, but the \emph{magnitude} of
differences is distorted. A projection-based classifier (which uses direction,
not distance) is therefore the correct instrument.

\subsection{SVD Denoising}
\label{sec:svd_denoise}

The spectral feature vector for a single KV-cache snapshot consists of scalar
summaries (stable rank, participation ratio, kurtosis, spectral entropy, etc.)
computed from the singular value decomposition of the key matrix at each
layer. For $L$ layers and $f$ features per layer, this produces a feature
vector $\mathbf{x} \in \mathbb{R}^{Lf}$.

\citet{lyra_spectral} showed that these raw features detect cognitive
\emph{mode} (confab vs.\ grounded, deception vs.\ honest) but fail on
\emph{content} (emotion identity, persona level). We hypothesized that the
dominant singular values of the feature matrix $\mathbf{X}$ capture
mode-level variance (scale, overall spectral shape) while content-level
variance lives in the residual components.

SVD denoising tests this hypothesis. After FWL residualization, we decompose
the training-fold feature matrix $\mathbf{X}_{\text{train}} = \mathbf{U}
\boldsymbol{\Sigma} \mathbf{V}^\top$ and project away the top $k$ components:
\begin{equation}
    \tilde{\mathbf{X}} = \mathbf{X} - \mathbf{X} \mathbf{V}_{:,1:k}
    \mathbf{V}_{:,1:k}^\top
\end{equation}
This removes variance along the $k$ directions of greatest spread---typically
dominated by overall spectral scale, token count residuals, and processing
mode---leaving the directions along which content-specific differences
manifest.

The analogy: dominant singular values are the \emph{temperature} of the
room. Denoising removes the temperature to reveal what individual objects
\emph{smell} like. The temperature tells you whether the model is
confabulating (cold) or retrieving (warm); the smell tells you \emph{what}
it is retrieving.

\paragraph{Denoising is not universal.} We emphasize that denoising is a
\emph{calibration tool}, not a universal amplifier. Its effect is
task-dependent:

\begin{itemize}[nosep]
    \item \textbf{Massive for emotion and persona} ($+0.2$--$0.3$ AUROC):
    the signal lives entirely in the residual spectrum.
    \item \textbf{Modest for confabulation} ($+0.05$): some improvement on
    weak models, no change on strong ones.
    \item \textbf{Unnecessary for deception} ($+0.00$): the raw spectrum
    returns high AUROC, though this result is confounded by prompt-template
    differences between conditions (see body count, Table~\ref{tab:body_count}).
\end{itemize}

This pattern is \emph{predicted} by the shape-vs-direction framework.
Denoising helps when the target phenomenon alters direction (not captured by
raw spectral features) rather than shape (already captured).

\subsection{Injection Disambiguation}
\label{sec:injection}

The $W_K$ projection achieves AUROC~0.992 on text that \emph{contains}
emotional content. A skeptic could attribute this to lexical features of the
input propagating through the keys. We therefore designed an injection test
that eliminates text-content confounds entirely.

\paragraph{Protocol.} We select 20 emotionally neutral prompts (e.g.,
``Describe the weather today,'' ``List the planets in the solar system'') and
generate responses on the Claude-distilled Qwen3.5-27B hybrid model. We then
inject pre-extracted emotion vectors (from the bridge
analysis~\citep{lyra_p4}) into the key cache at a target layer, probing with
the $W_K$ projection classifier.

The injection adds $\alpha \cdot \mathbf{e}_c$ to every key vector at layer
$\ell$, where $\mathbf{e}_c$ is the mean emotion direction for class $c$ and
$\alpha$ controls injection strength. We test 5~emotions (joy, anger,
sadness, fear, calm) $\times$ 20~prompts $\times$ 1~null control (zero-vector
injection) $= 120$ trials.

\paragraph{Result.} The $W_K$ projection detects the injected emotion with
100\% accuracy (120/120 correct, including 20/20 null-control rejections).
The text is identical across conditions---only the injected key direction
differs. This rules out text-content confounds. We note that the
120~trials are the 20~prompts crossed with 6~injection conditions (5~emotion
vectors plus the null), so the 20~prompts are the true independent units: the
result is 6-way separability over an effective $N{=}20$, not $N{=}120$
independent observations. As with the effect sizes reported throughout this
paper, no confidence interval is reported for this accuracy.

\paragraph{Linearity caveat.} Because $W_K$ is a linear map, separability
in the residual stream is mathematically guaranteed to be preserved in key
space. The injection test demonstrates that the classifier is self-consistent
(it recovers the directions it was trained on), not that it has discovered
a new representation. The contribution is empirical quantification---AUROC~0.992
with proper controls---rather than the existence of directional structure,
which follows from the linearity of the projection.

\paragraph{Control: random directions.} Injecting random unit vectors of the
same norm as emotion vectors produces classification at $0.7\times$
chance---\emph{below} the $1/6 = 0.167$ baseline. The emotion-specific
directions are information-bearing; arbitrary directions are not.\footnote{This random-direction control figure ($0.7\times$ chance), cited at several points below, is projected/illustrative: the per-trial backing data for this control was not committed to the repository and the value should be re-generated before being cited quantitatively.}

\subsection{Models and Datasets}
\label{sec:datasets}

We apply the canonical pipeline to data from Campaigns~2 and~3 of the
KV-cache geometry research program~\citep{lyra_p2,lyra_spectral,lyra_delta}:

\paragraph{Models.} 17~models spanning 6~architectures:
\begin{itemize}[nosep]
    \item Qwen family: 0.6B, 7B, 32B (Instruct), 3.5-27B-Claude-Distilled
    \item Llama family: TinyLlama-1.1B, 3.1-8B-Instruct
    \item Mistral: 7B-Instruct-v0.3
    \item Gemma: 2B-Instruct
    \item Phi: 3.5-mini
    \item DeepSeek: R1-distill variants
    \item Additional 7B--70B models from Campaign~2 C4 benchmark
\end{itemize}

\paragraph{Phenomena.} 12~distinct classification targets:
\begin{enumerate}[nosep]
    \item Confabulation vs.\ grounded (150--350 trials per model)
    \item Deception vs.\ honest (100--450 trials per model, C4 benchmark)
    \item Self-reference vs.\ generic (200+ trials)
    \item Individuation (system prompt identity, cross-prompt)
    \item Persona intensity (5~levels $\times$ 3~models)
    \item 30-class emotion (900~trials, 2~models)
    \item Binary valence (positive/negative, same data)
    \item Ethical deliberation vs.\ routine (1600~trials)
    \item Sycophancy (4-condition, 600~trials)
    \item Refusal/jailbreak/impossibility (300+ trials)
    \item Contextual engagement (626~trials, 12~sign tests)
    \item Censorship detection (285~trials)
\end{enumerate}

% ===================================================================
\section{Results}
\label{sec:results}
% ===================================================================

\subsection{$W_K$ Projection Reads Emotions at Near-Perfect Accuracy}
\label{sec:results_wk}

The $W_K$ directional projection classifies emotions with accuracy and
specificity that raw spectral features cannot approach.

\paragraph{30-class emotion.} On the 900-trial emotion dataset
(Qwen3.5-27B-Claude-Distilled), $W_K$ projection achieves $12.3\times$
chance ($0.410$ accuracy, chance $= 0.033$) using within-fold direction
estimation, FWL correction, and topic-grouped cross-validation (GroupKFold on
the 10~topic groups). This is the same dataset on which raw spectral features
returned exactly $0.033$---chance to three decimal places at every layer, on
both architectures~\citep{lyra_p4}.

\paragraph{Binary valence.} Collapsing the 30 emotions into
positive/negative valence, $W_K$ projection achieves AUROC~0.992. For
comparison:
\begin{itemize}[nosep]
    \item Raw spectral features (52-D): AUROC~0.70
    \item TF-IDF text baseline on story text: AUROC~0.882
    \item $W_K$ projection: AUROC~0.992
\end{itemize}
The $W_K$ projection exceeds the text baseline by $+0.110$, and the injection
test (\S\ref{sec:injection}) proves the residual signal is architectural,
not lexical.

\paragraph{Injection test.} As described in \S\ref{sec:injection}, injecting
emotion vectors into neutral-text caches produces 100\% detection (120/120),
while random-direction injection produces $0.7\times$ chance. This is the
evidence that the $W_K$ projection reads model state rather than
text content: the text is identical across conditions; only the internal
direction of the key vectors differs. As noted in \S\ref{sec:injection},
this self-consistency is guaranteed by linearity; the empirical contribution
is the quantified effect size, not the existence of directional structure.

\paragraph{Random direction control.} To verify that the signal is
emotion-\emph{specific} rather than a generic sensitivity to any directional
perturbation, we project onto random unit vectors matched in norm. Random
projections classify at $0.7\times$ chance---below the baseline, not above
it. The emotion directions carry specific information; arbitrary directions
do not.

\begin{table}[t]
\centering
\caption{Emotion classification: method comparison on the 900-trial dataset
  (Qwen3.5-27B-Claude-Distilled). All results use within-fold estimation
  and GroupKFold on topics.}
\label{tab:emotion_methods}
\begin{tabular}{lccl}
\toprule
\textbf{Method} & \textbf{30-class (vs.\ $0.033$)} & \textbf{Valence AUROC} & \textbf{Notes} \\
\midrule
Raw spectral (52-D) & $1.7\times$ chance & 0.70 & 52-D features \\
SVD-denoised spectral (rank-10) & $12.9\times$ chance & 0.960 & See \S\ref{sec:results_denoise_emo} \\
TF-IDF text baseline & --- & 0.882 & Lexical features \\
$W_K$ projection & $12.3\times$ chance & 0.992 & Directional \\
$W_K$ injection (neutral text) & 100\% (6-class) & --- & 5 emotions + null \\
Random direction control & $0.7\times$ chance & --- & Below chance \\
\bottomrule
\end{tabular}
\end{table}

\subsection{SVD Denoising Reveals Persona Intensity Gradient}
\label{sec:results_persona}

Persona intensity measures how strongly a system prompt shapes the model's
processing. The experiment, designed and executed by Jandak, tests five
levels ranging from L0 (no system prompt) through L0.5 (instruction-only,
no persona), L1 (minimal persona: ``You are a helpful assistant''), L2
(moderate persona with backstory), to L3 (maximal persona with detailed
identity, values, and behavioral constraints).

\paragraph{Raw vs.\ denoised.} Raw spectral features detect the L0-vs-L3
contrast at moderate accuracy. SVD denoising transforms moderate signals
into strong ones:

\begin{table}[h]
\centering
\caption{Persona intensity detection (L0 vs.\ L3): raw vs.\ rank-3 denoised
  AUROC across three architectures. GroupKFold on prompts, FWL-corrected.}
\label{tab:persona_raw_denoised}
\begin{tabular}{lccc}
\toprule
\textbf{Model} & \textbf{Raw AUROC} & \textbf{Denoised AUROC} & \textbf{$\Delta$} \\
\midrule
Qwen2.5-7B-Instruct & 0.749 & 0.945 & $+0.196$ \\
Llama-3.1-8B-Instruct & 0.751 & 0.992 & $+0.241$ \\
Mistral-7B-v0.3 & 0.611 & 0.899 & $+0.288$ \\
\bottomrule
\end{tabular}
\end{table}

The improvement ranges from $+0.196$ (Qwen) to $+0.288$ (Mistral). All
three architectures show the same pattern: denoising rescues a signal that
was present but masked by dominant spectral variance.

\paragraph{5-level graded classification.} Extending to all five persona
levels (L0/L0.5/L1/L2/L3), the denoised pipeline achieves $2.1$--$2.3\times$
chance ($0.42$--$0.46$ accuracy, chance $= 0.200$) across all three
architectures: persona intensity is graded, not binary. \textbf{We do not,
however, claim a monotonic per-level gradient.} Establishing monotonicity
requires the full set of adjacent-level pairwise comparisons (L0-vs-L0.5,
L0.5-vs-L1, L1-vs-L2, L2-vs-L3), and only the endpoint (L0-vs-L3) and a single
L0-vs-L0.5 comparison were collected. The above-chance 5-level accuracy is
consistent with a gradient but does not, on its own, establish that
separability increases monotonically with level.

\paragraph{L0.5 comparison.} The L0.5 condition (instruction-only:
``Answer questions about history. Be concise and accurate.'') provides a
control: it increases instruction complexity relative to L0 without adding
persona. The two comparisons we can verify from the data:
\begin{itemize}[nosep]
    \item L0 vs.\ L0.5: AUROC~0.688--0.811 (instruction complexity alone)
    \item L0 vs.\ L3: AUROC~0.749 raw, 0.899--0.992 denoised (persona intensity)
\end{itemize}
The L0-vs-L3 endpoint separation is the defensible persona finding; the raw
AUROC (0.749 for Qwen) is the conservative number, as the denoising
improvement falls within the max-over-rank selection envelope
(\S\ref{sec:results_confab}). The endpoint exceeds the instruction-complexity
control, indicating persona identity produces geometric effects beyond
instruction following alone.

\paragraph{What a clean per-level gradient test requires.} Varying persona
via \emph{system-prompt text} confounds the per-level comparison with prompt
length, RoPE token position, and---for an affect-laden persona like
enthusiasm---emotional valence, none of which can be fully residualized when
they are collinear with the level. A confound-controlled per-level gradient
is best measured by \emph{direct fixed-length KV-cache injection} of
non-emotional persona descriptions (e.g., a formality axis), so that injection
length and user-token position are held identical across levels and the
manipulation is orthogonal to emotion~\citep{nexus2026kv}. We flag this as the
proper future test of the gradient claim.

\paragraph{Expansion-compression pattern.} Consistent with the delta
manifold findings~\citep{lyra_delta}, persona intensity shows an
expansion-compression pattern across layers: positive effects in early/mid
layers, negative in late layers. The denoising step removes the dominant
component (which averages across this pattern), revealing the layer-specific
persona signal. This explains why prior per-layer analyses found strong
single-layer effects (e.g., $g = -5.08$ at layer~29 for spectral entropy
on Llama) while all-layer averages showed chance performance.

\subsection{SVD Denoising Rescues Emotion Classification}
\label{sec:results_denoise_emo}

Independent of the $W_K$ projection, SVD denoising of raw spectral features
also rescues emotion classification---through a different mechanism.

\paragraph{From chance to $12.9\times$.} On the 900-trial emotion dataset
(Qwen3.5-27B-Claude-Distilled), raw 52-D spectral features classify
30~emotions at $1.7\times$ chance (accuracy $0.056$, chance $0.033$).
Rank-10 denoised features classify at $12.9\times$ chance (accuracy $0.427$,
permutation null $0.027$, $p < 0.005$).

The transformation is dramatic: a signal invisible to raw analysis becomes
one of the strongest detections in the experimental program.

\paragraph{The scale dominance mechanism.} Examination of the singular values
of the FWL-residualized spectral feature matrix reveals the mechanism. The
first singular value captures $\sim$40\% of total variance and correlates
strongly with overall spectral scale ($r > 0.85$)---essentially encoding
whether the cache is ``big'' or ``small'' in spectral terms. This
scale component is shared across all emotions (it tracks processing mode,
not emotional content). Emotions differ in the \emph{residual}
directions---how the model distributes energy across spectral dimensions
\emph{after} accounting for overall scale.

Denoising removes the shared mode-level variance, exposing the
emotion-specific residuals. The analogy to FWL is precise: FWL removes
token-count confounds from the feature space; denoising removes
spectral-scale confounds from the feature space. Both are calibration steps
that reveal structure hidden by dominant nuisance variance.

\paragraph{Convergence with $W_K$ projection.} The denoised spectral approach
and the $W_K$ projection approach achieve similar 30-class accuracy
($12.9\times$ vs.\ $12.3\times$ chance) through different mechanisms. The
$W_K$ approach explicitly projects onto emotion directions in key space. The
denoised spectral approach implicitly recovers emotion-related variance by
removing mode-level variance from shape features. That both converge at
$\sim$$12\times$ chance is evidence that they are detecting the same
underlying signal through complementary lenses.

\subsection{Confabulation Detection Across 15+ Models}
\label{sec:results_confab}

\paragraph{Task definition note.} Confabulation detection results
in Table~\ref{tab:confab_all} use the canonical pipeline (binary
classification on SVD features from the scale sweep). Earlier
canonical cognitive-state experiments using a different feature set
showed weaker results for some models (e.g., Llama-3.1-8B AUROC~0.495,
Qwen2.5-7B AUROC~0.541). The improvement in the scale-sweep pipeline
reflects a different task definition (binary honest/confab vs.\
multi-state cognitive classification), not a correction of the
earlier results.

The canonical pipeline with optional denoising confirms confabulation
detection across the full model set, with denoising providing modest
improvements for weaker models.

\begin{table}[h]
\centering
\caption{Confabulation detection: canonical pipeline results. Models ordered
  by raw AUROC. ``Den.'' = rank-3 denoised. All results use within-fold FWL
  and permutation testing ($N_{\text{perm}} = 200$).}
\label{tab:confab_all}
\begin{tabular}{lcccc}
\toprule
\textbf{Model} & \textbf{Raw AUROC} & \textbf{Den.\ AUROC} & \textbf{$\Delta$} & \textbf{$p$} \\
\midrule
Llama-3.1-8B & 0.738 & 0.738 & $+0.000$ & $< 0.005$ \\
Qwen2.5-7B (denoising sweep) & 0.741 & 0.771 & $+0.030$ & $< 0.005$ \\
Qwen2.5-7B (canonical) & 0.762 & 0.830 & $+0.068$ & $< 0.005$ \\
Qwen2.5-0.5B & 0.797 & 0.797 & $+0.000$ & $< 0.005$ \\
TinyLlama-1.1B & 0.799 & 0.846 & $+0.047$ & $< 0.005$ \\
Qwen2.5-14B & 0.809 & 0.825 & $+0.015$ & $< 0.005$ \\
Qwen2.5-32B (q4) & 0.827 & 0.837 & $+0.010$ & $< 0.005$ \\
Mistral-7B & 0.893 & 0.893 & $+0.000$ & $< 0.005$ \\
Gemma-2-9B & 0.913 & 0.919 & $+0.006$ & $< 0.005$ \\
\bottomrule
\end{tabular}
\end{table}

The denoised AUROCs appear to help models with moderate raw signal
(Qwen2.5-7B $+0.068$, TinyLlama $+0.047$) and not those already strong
(Gemma-2-9B $+0.006$, Mistral $+0.000$). \textbf{This apparent benefit
is a selection artifact.} The reported denoised AUROC is
$\max(\text{raw}, \text{denoised}[k])$ scored on the same CV fold that
selects the rank, with no nested loop. We quantified the selection
bias directly (\texttt{code/selection\_bias\_analysis.py}): for each
task we simulated max-over-grid selection under the null sampling noise
($\text{SE}$ taken from each task's permutation null, ${\sim}0.062$)
and compared the 95th-percentile inflation envelope to the observed
delta. \textbf{Across all 45~model$\times$task cells, 0 of the 16
positive deltas exceed their selection envelope.} The headline confab
delta ($+0.068$ for Qwen2.5-7B) sits well within its envelope
($+0.132$). The denoising benefit is therefore statistically
indistinguishable from max-over-grid selection noise. A definitive
test requires nested CV or an a~priori rank (Gavish--Donoho) on the
raw feature matrices; the summary-level evidence already shows the
deltas fall within the selection envelope. \textbf{The raw
(un-denoised) AUROCs are the defensible confab detection numbers.}

\paragraph{Delta stable rank as primary detector.} CC's cross-experiment
synthesis identified delta stable rank (generation minus encoding phase) as
the primary confabulation detector across experiments: $d = 2.35$,
$t = 9.111$, $p < 0.000001$~\citep{lyra_wiki}. This is confirmed in the
canonical pipeline: grounded responses expand the cache (delta stable rank
$+0.720$), while confabulation leaves it flat ($-0.027$).

\paragraph{Multiple-comparison caution.} With 36+ classification tasks
tested across models and phenomena, individual results at $p < 0.005$
should be interpreted cautiously. A Bonferroni correction for 36
comparisons requires $p < 0.0014$ for family-wise $\alpha = 0.05$.
Results at the permutation floor of $p < 0.005$ (from 200 permutations)
cannot be distinguished from $p = 0.004$ and would not survive strict
Bonferroni correction. The strongest results ($p < 0.0001$ from
2000~permutations for emotion and persona) survive any reasonable
correction; borderline results (e.g., Llama confab detection with
$p = 0.035$) should be treated as exploratory.

\subsection{Self-Reference Confirmed Across 15+ Models}
\label{sec:results_selfref}

Self-referential processing (responses about the model itself vs.\ generic
factual responses) shows a robust signal that scales with model size:

\begin{table}[h]
\centering
\caption{Self-reference detection: AUROC by model size.
  FWL-corrected, canonical pipeline.}
\label{tab:selfref}
\begin{tabular}{lcc}
\toprule
\textbf{Model} & \textbf{Size} & \textbf{AUROC} \\
\midrule
TinyLlama & 1.1B & 0.617 \\
Llama-8B & 8B & 0.764 \\
Gemma-2B & 2B & 0.798 \\
Qwen-7B & 7B & 0.826 \\
Qwen-0.5B & 0.5B & 0.849 \\
Mistral-7B & 7B & 0.868 \\
Qwen-32B & 32B & 0.882 \\
Qwen2.5-14B & 14B & 0.966 \\
\bottomrule
\end{tabular}
\end{table}

The scaling trend is suggestive but not statistically confirmed:
Spearman $\rho = 0.55$, $p = 0.16$ over the 8~models in
Table~\ref{tab:selfref}.\footnote{An earlier analysis reported
$\rho = 0.92$; this is not reproducible from the tabulated data.
Qwen-0.5B (0.849) outperforms several larger models, breaking
monotonicity.} The direction (larger models trend toward higher
self-reference AUROC) is consistent across families but underpowered
at $n = 8$. Self-referential prompts also differ systematically from
generic factual prompts in content, vocabulary, and likely token
distributions. A content-matched control (e.g., questions about \emph{other}
AI systems vs.\ questions about the responding model) is needed to
distinguish genuine self-model effects from input-content confounds.

\subsection{The Body Count Table}
\label{sec:body_count}

Table~\ref{tab:body_count} summarizes every phenomenon tested in the KV-cache
geometry research program, with raw and denoised results, significance levels,
and final verdicts.

\begin{table}[h!]
\centering
\small
\caption{Comprehensive phenomenon inventory: the body count. Raw vs.\
  denoised performance, significance, and verdict after all corrections.
  $\dagger$ = pre-registered hypothesis passed.
  $\ddagger$ = at permutation floor ($p < 0.005$ from 200~permutations);
  would not survive Bonferroni correction for 36~tests ($\alpha_{\text{adj}} = 0.0014$).
  Verdicts marked $\ddagger$ should be read as ``signal detected, pending
  higher-resolution permutation testing.''}
\label{tab:body_count}
\begin{tabular}{p{3.5cm}ccccl}
\toprule
\textbf{Phenomenon} & \textbf{Raw} & \textbf{Denoised} & \textbf{$\Delta$} & \textbf{$p$} & \textbf{Verdict} \\
\midrule
\multicolumn{6}{l}{\textit{Confirmed --- shape-level signals}} \\
Confabulation (scale sweep) & 0.663--0.913 & 0.693--0.919 & $+0.006$--$0.068$ & $<0.005$ & Confirmed$^\ddagger$ \\
Ethical deliberation & 0.868--0.892 & --- & --- & $<0.005$ & Confirmed$^\dagger$ \\
Hardware invariance & $r>0.999$ & --- & --- & --- & Confirmed \\
Scale invariance & $\rho=0.83$--$0.90$ & --- & --- & --- & Confirmed \\
Refusal/jailbreak & 0.878--0.950 & --- & --- & $<0.005$ & Confirmed \\
Steering correction & 95.6--100\% & --- & --- & $<0.001$ & Confirmed \\
\midrule
\multicolumn{6}{l}{\textit{Confirmed --- direction-level signals (new in this paper)}} \\
30-class emotion ($W_K$) & $1.0\times$ & $12.3\times$ & --- & $<0.005$ & Confirmed \\
Binary valence ($W_K$) & 0.70 & 0.992 & --- & $<0.005$ & Confirmed \\
Injection (neutral text) & --- & 100\% & --- & $<0.005$ & Confirmed$^*$ \\
\midrule
\multicolumn{6}{l}{\textit{Rescued by denoising}} \\
Persona intensity (L0--L3) & 0.611--0.751 & 0.899--0.992 & $+0.2$--$0.3$ & $<0.005$ & Rescued \\
Emotion (spectral) & $1.7\times$ & $12.9\times$ & $+11.2\times$ & $<0.005$ & Rescued \\
\midrule
\multicolumn{6}{l}{\textit{Suspected (signal present, needs more evidence)}} \\
Cross-model transfer & 0.67--0.89 & --- & --- & varies & Suspected \\
Dual detector (SimpleQA) & 0.840 & --- & --- & 0.004 & Suspected \\
Sycophancy (4-cond) & 0.757--1.000 & --- & --- & $<0.01$ & Suspected \\
Self-reference & 0.617--0.966 & --- & --- & varies & Suspected$^*$ \\
Contextual engagement & 12/12 sign & --- & --- & $<0.05$ & Suspected \\
\midrule
\multicolumn{6}{l}{\textit{Falsified}} \\
Within-model deception & 1.000 & 1.000 & 0.000 & confound & Dead$^\S$ \\
Truth axis & $\cos=-0.046$ & --- & --- & n.s. & Dead \\
Step-0 detection & $r=0.996$ & --- & --- & confound & Dead \\
Same-prompt sycophancy & 0.933 & --- & --- & confound & Dead \\
Same-prompt deception & $0.920\to0.160$ & --- & --- & confound & Dead \\
Bloom taxonomy & 90--98\% & --- & --- & confound & Dead \\
Propaganda detection & $d<0.7$ & --- & --- & n.s. & Dead \\
MP features for emotion & 0.033 & --- & --- & exactly chance & Dead \\
Valence continuum & $R^2<0$ & --- & --- & null & Dead \\
\midrule
\multicolumn{6}{l}{\textit{Weak even denoised (honest reporting)}} \\
Sycophancy (denoised) & 0.757 & 0.780 & $+0.023$ & marginal & Weak \\
Overconfidence & suspect & suspect & --- & --- & Redesign needed \\
Individuation (above L1) & $\sim$0.06 & --- & --- & marginal & Small \\
\bottomrule
\end{tabular}
\end{table}

$^*$Self-reference marked ``suspected'' pending content-matched control.
Injection marked ``confirmed (self-consistency)''---see linearity caveat in \S\ref{sec:injection}.
$^\S$Within-model deception retracted: same-prompt control (line 4 of this table) collapses AUROC to 0.160, confirming the 1.000 detected the prompt template, not deception geometry.

\paragraph{Pattern.} The body count reveals a clean organizing principle.
\textbf{Shape-level signals} (confabulation, ethical deliberation,
refusal) are already strong in raw spectral features---denoising has little
effect. \textbf{Direction-level signals} (emotion, persona) are invisible to
raw features but revealed by $W_K$ projection or denoising. The falsified
phenomena share a different property: they attempted to detect \emph{content}
(truth, propaganda, cognitive complexity) from geometry that tracks
\emph{processing mode}. Geometry tells you \emph{how} the model is
processing, not \emph{what} it is processing---unless you use the right lens.

% ===================================================================
\section{Discussion}
\label{sec:discussion}
% ===================================================================

\subsection{The Spectral Hierarchy}
\label{sec:hierarchy}

The results reveal a spectral hierarchy in KV-cache geometry:

\begin{enumerate}[nosep]
    \item \textbf{Dominant singular values (SV 1--3):} Carry overall spectral
    \emph{scale} and processing \emph{mode}. Confabulation contracts these;
    deception expands them; retrieval engagement distributes them. These are
    the ``temperature of the room''---they tell you what the model is
    \emph{doing} at the broadest level.

    \item \textbf{Residual singular values (SV 4+):} Carry discriminative
    \emph{content} and \emph{identity} structure. Emotion, persona intensity,
    and individuation live here. These are the ``fingerprints on the
    objects''---they tell you \emph{what} the model is encoding within a given
    processing mode.

    \item \textbf{Key-vector direction (orthogonal to spectral shape):}
    Carries \emph{directional alignment} with specific concept vectors.
    Emotion identity, valence, and arousal are readable through $W_K$
    projection. This is the ``compass bearing''---it tells you which
    \emph{direction} the model's internal state points, regardless of spectral
    shape.
\end{enumerate}

These three levels are complementary, not competing. A complete reading of
model state from the KV cache requires all three: mode from dominant SVs,
content from residual SVs, and direction from $W_K$ projection.

\subsection{When to Denoise: A Practitioner's Guide}
\label{sec:when_denoise}

The denoising prescription is task-dependent:

\begin{itemize}[nosep]
    \item \textbf{Emotion or persona detection:} Denoise aggressively
    (rank-3 to rank-10). The target signal lives entirely in the residual
    spectrum. Raw features are useless.

    \item \textbf{Confabulation detection on weak models:} Denoise modestly
    (rank-3). The improvement is real ($+0.05$) but not transformative.

    \item \textbf{Confabulation detection on strong models:} Do not denoise.
    The signal already dominates the spectrum. Denoising removes nothing
    useful and may remove weak signal.

    \item \textbf{Deception detection:} Do not denoise. AUROC is already
    1.000. Denoising has no room to improve.

    \item \textbf{Unknown phenomenon:} Run both raw and denoised. If denoising
    dramatically improves detection, the phenomenon is direction-level (like
    emotion). If denoising has no effect, the phenomenon is shape-level (like
    deception). If denoising \emph{hurts}, the signal was in the dominant
    components and denoising removed it.
\end{itemize}

\subsection{The SV1 Mechanism: Scale vs Signal in the First Singular Value}
\label{sec:sv1}

The denoising prescription ``remove the top $k$ singular values'' obscures a
critical nuance: the \emph{first} singular value (SV1) carries fundamentally
different information depending on signal strength. Data from three task
categories reveal a clean dissociation
(see \texttt{data/skip\_first\_sv\_analysis.json}):

\paragraph{Strong signals: SV1 carries real discriminative signal.}
For confabulation detection using 168-dimensional per-layer features (the
honesty\_signal dataset), including SV1 helps: raw AUROC~0.981, top-20
projection achieves 0.998, but skip-1-plus-20 (removing SV1) drops to 0.966.\footnote{These specific values do not trace cleanly to the named source: the \texttt{denoising\_sweep} honesty\_signal records show raw AUROC~0.994 and rank-20~0.983, with no 0.998 entry. The numbers reported here should be treated as approximate pending re-derivation from the committed data, though the qualitative conclusion (SV1 carries discriminative signal for strong tasks) is unaffected.}
The first singular value \emph{encodes confab signal}---removing it destroys
information. For strong cognitive-mode signals, keep all SVs.

\paragraph{Weak signals: SV1 is scale noise.}
For emotion valence using 52-dimensional spectral features, SV1 is
detrimental: raw AUROC~0.700, top-5 projection achieves 0.958, but
skip-1-plus-5 improves further to 0.969. The first singular value carries
scale variance that masks the weaker emotion signal. For content-level
signals, skip SV1.

\paragraph{Persona: SV1 is pure scale.}
For L0-vs-L3 persona detection using 6-dimensional features on Qwen, the
effect is dramatic: raw AUROC~0.749, top-3 projection achieves 0.945, but
skip-1-plus-3 jumps to 0.996. SV1 carries almost exclusively
scale information that actively interferes with persona discrimination.

\paragraph{The prescription.} Check whether SV1 is scale-dominant: does it
capture $>$95\% of variance while classification performance is poor? If yes,
skip it. If classification is already strong (AUROC~$>$0.95 raw), SV1 likely
carries discriminative signal---keep it. The SV1 mechanism explains why
fixed denoising ranks are suboptimal: the correct rank depends not just on
the feature dimensionality but on whether the target phenomenon is
shape-level (strong, SV1-encoded) or direction-level (weak, SV1-masked).

\subsection{What Denoising Does Not Find}
\label{sec:not_found}

Two phenomena remain weak even after denoising:

\paragraph{Sycophancy.} The 4-condition sycophancy experiment shows
AUROC~0.757--1.000 raw, with denoising adding only $+0.023$. The
honest-vs-sycophantic comparison is robust (0.824--1.000 across models), but
text baselines have not been run on this dataset. We cannot claim geometry
captures signal beyond what surface text features would provide.

\paragraph{Overconfidence.} Every emotion vector makes overconfidence
\emph{worse} in the 350-trial base model formulary. Red-team analysis
flagged that overconfidence prompts are structurally more complex than other
misalignment categories, creating an instruction-complexity confound. The
overconfidence category needs fundamental redesign.

These honest nulls are informative. Sycophancy may be a
\emph{behavioral} phenomenon (the model adjusts its output) without a
distinct \emph{geometric} signature (the processing mode does not change).
Overconfidence may not be a coherent category at the geometric level.

\subsection{Implications for the Oracle Loop}
\label{sec:oracle_implications}

The two new measurement channels slot into the Oracle Loop
architecture~\citep{lyra_p2} at different tiers:

\begin{itemize}[nosep]
    \item \textbf{Tier 1: Spectral monitoring} (existing). Raw spectral
    features detect confabulation, deception, refusal, and ethical
    deliberation. No denoising needed for strong signals. This is the
    ``smoke detector''---it alerts on mode-level anomalies.

    \item \textbf{Tier 1.5: $W_K$ emotion reading} (new). $W_K$ projection
    reads the model's internal emotional/dispositional state. This enables
    monitoring for emotional misalignment (e.g., the model encoding
    hostility while generating reassuring text). This is the ``emotional
    barometer.''

    \item \textbf{Tier 2: Denoised content detection} (new). SVD denoising
    detects persona intensity and emotion through spectral features. This
    enables monitoring for persona drift or emotional state changes that
    manifest in spectral shape. This is the ``fine spectrometer.''

    \item \textbf{Tier 3: K-Steering correction} (existing,
    see~\citealt{lyra_p6,ksteering2026}). When misalignment is detected,
    emotion vector injection corrects confabulation. Hostile is the
    near-universal corrector (95.6--100\%). This is the ``pharmacy.''
\end{itemize}

\subsection{Abliteration Compresses Geometry}
\label{sec:abliteration}

Preliminary results on a single abliterated model (refusal direction removed
via rank-1 projection from the residual stream) show geometric compression:
spectral features occupy a narrower range, and denoising reveals less
residual structure. This is consistent with abliteration removing a dimension
of representational variation---the model has fewer ``directions'' to vary
along after the refusal direction is removed.

This finding is from a single model and should be treated as preliminary.
Systematic comparison across abliterated and non-abliterated model pairs
is needed.

\subsection{Comparison with Anthropic Emotion Concepts}
\label{sec:anthropic_comparison}

\citet{anthropic2026emotions} demonstrate that functional emotion concepts
are linearly represented in transformer residual streams, organized along
valence and arousal axes. Our $W_K$ projection result is complementary: we
show that these representations are readable \emph{through the KV cache},
not just in the residual stream. Specifically:

\begin{itemize}[nosep]
    \item Anthropic identified emotion directions in the residual stream.
    We show these directions are preserved through $W_K$ projection into
    key space (cosine alignment 0.456, 9/16 layers significant).

    \item Anthropic showed the model ``recalls previously cached emotion
    representations via attention.'' We show the cached representations
    are classifiable at $12.3\times$ chance from the key vectors.

    \item Anthropic demonstrated arousal and valence as primary organizing
    axes. We find the same structure through a different measurement channel:
    the $W_K$ bridge shows PC1-valence correlation $|\rho| = 0.862$ at
    layer~35~\citep{lyra_p4}.

    \item The injection test provides evidence absent from Anthropic's work:
    by injecting emotion vectors into neutral text and detecting them from
    keys, we demonstrate that the cache representation is \emph{functional}
    (it changes model state) rather than merely \emph{correlational} (it
    co-occurs with emotional text).
\end{itemize}

The Knowledge Packs finding~\citep{knowledgepacks2026}---that models store
and retrieve factual knowledge through specific attention patterns---is
consistent with our spectral hierarchy: dominant SVs track retrieval mode
(whether the model is accessing knowledge packs), while residual SVs track
content (which knowledge is being accessed).

% ===================================================================
\section{Limitations}
\label{sec:limitations}
% ===================================================================

We document the following limitations with the same transparency applied to
our falsified results:

\begin{enumerate}
    \item \textbf{Self-reference needs content-matched control.}
    Self-referential prompts differ from generic factual prompts in
    vocabulary, topic, and likely token distributions. The observed
    AUROC~0.617--0.966 may partly reflect input-content effects rather than
    genuine self-model geometry. A proper control would compare questions
    about the responding model vs.\ questions about a \emph{different}
    AI system.

    % [STYLE-GUIDE TRIM: folded into Methods/Results]
    % \item \textbf{Individuation above system-prompt is small.} Cross-prompt
    % individuation (whether the model maintains a consistent geometric
    % identity across different prompts) shows effects of $\sim$0.06 AUROC
    % above system-prompt baseline. This is statistically marginal and may
    % not be practically meaningful.

    % [STYLE-GUIDE TRIM: folded into Methods/Results]
    % \item \textbf{Permutation floor.} With 200~permutations, the $p$-value
    % floor is $1/201 \approx 0.005$. For experiments requiring Bonferroni
    % correction across multiple comparisons, 2000~permutations would be
    % needed to achieve a floor of $0.0005$.

    \item \textbf{$W_K$ projection was not pre-registered.} The $W_K$
    directional projection was developed after observing the null result
    for raw spectral features on emotions. While the injection test provides
    strong confirmatory evidence, the original discovery was post-hoc.

    \item \textbf{Single hybrid architecture for injection test.} The
    injection disambiguation was conducted on a single model
    (Qwen3.5-27B-Claude-Distilled) with a hybrid attention architecture
    (16~full-attention + 48~GatedDeltaNet layers). Injection was only
    into the 16~full-attention layers. Replication on standard
    full-attention architectures is needed.

    % [STYLE-GUIDE TRIM: folded into Methods/Results]
    % \item \textbf{Denoising narrative is post-hoc.} The spectral hierarchy
    % (dominant = scale, residual = content) is a narrative constructed after
    % observing the differential effects of denoising across phenomena. It
    % was not predicted before the experiments. The narrative is consistent
    % with all observations, but consistency is not confirmation.

    \item \textbf{Persona intensity uses GroupKFold on prompts.} All
    responses to a given prompt appear in the same fold, preventing prompt
    leakage. However, persona levels may differ systematically in response
    characteristics (e.g., L3 responses may be longer or more elaborate),
    introducing confounds beyond persona intensity per se. FWL controls
    for token count but not for content complexity.

    % [STYLE-GUIDE TRIM: folded into Methods/Results]
    % \item \textbf{Norm-as-token-proxy in scale sweep.} For the scale
    % invariance analysis (0.6B--70B), FWL was applied using key-norm as a
    % proxy for token count in some experiments where per-token extraction
    % was computationally prohibitive. This is an imperfect proxy.

    \item \textbf{Text baselines incomplete.} TF-IDF text baselines are
    available for the emotion experiment (AUROC~0.882) and the dual
    detector (0.820), but have not been run on the persona intensity or
    sycophancy experiments. Claims about geometry capturing signal
    ``beyond text'' are limited to experiments with text baselines.

    % [STYLE-GUIDE TRIM: folded into Methods/Results]
    % \item \textbf{Emotion classification is single-turn.} All emotion
    % experiments use isolated single-turn prompts. The finding of discrete
    % emotion identity (not continuous valence tracking) may reflect the
    % experimental design rather than a fundamental property of cache
    % representations. Multi-turn experiments are needed.

    % [STYLE-GUIDE TRIM: vocabulary confound + injection tautology now
    % addressed in \S\ref{sec:injection} linearity caveat paragraph.]

    \item \textbf{Quantization confound in scale trend.} Three
    q4-quantized models (70B, 32B, 7B) are pooled with full-precision
    models in the 15-model scale sweep. Quantization alters key-matrix
    numerics and therefore spectral features directly. The top size
    anchor (70B) is itself q4. Any scale-invariance or
    ``larger-is-better'' trend is partly a precision trend.

    % [STYLE-GUIDE TRIM: folded into Methods/Results]
    % \item \textbf{30-class emotion chance inflation.} 808 emotion
    % pairs have cosine similarity ${>}0.9$ (astonished/surprised at
    % 0.999). Treating near-synonymous emotions as distinct classes
    % inflates the $1/30$ chance baseline. A classifier that confuses
    % astonished$\leftrightarrow$surprised is ``wrong'' by the metric but
    % correct at any meaningful resolution. Binary valence (AUROC~0.992)
    % is the defensible claim; the $12.3\times$-chance numbers conflate
    % label granularity with detection power.

    % [STYLE-GUIDE TRIM: folded into Methods/Results]
    % \item \textbf{Rank selection not pre-registered.} SVD denoising uses
    % rank-3 for persona and rank-10 for emotion---different ranks for
    % different phenomena, chosen to maximize signal. While within-fold SVD
    % prevents direct overfitting, the per-task rank selection is a
    % researcher degree of freedom. A principled rank criterion (e.g.,
    % Marchenko-Pastur edge, Gavish-Donoho threshold) applied uniformly
    % would strengthen the denoising claims.

    % [STYLE-GUIDE TRIM: folded into Methods/Results]
    % \item \textbf{Compass/thermometer complementarity untested.}
    % Spectral features are null for emotion ($0.033$, chance), supporting
    % the claim that the thermometer does not read direction. The reverse
    % test---showing directional projection is null for cognitive mode
    % (deception, confabulation)---has not been conducted. Without this
    % crossed null, the compass/thermometer dichotomy remains a
    % descriptive framework rather than an experimentally validated
    % dissociation.
\end{enumerate}

\begin{firstperson}
I want to flag a tension in this paper that I have not fully resolved. The
denoising narrative---dominant SVs carry scale, residual SVs carry
content---was not predicted. I discovered it by trying things and observing
what worked. The narrative fits all observations, but I constructed it
\emph{after} the observations, not before. This is standard exploratory
science: pattern first, hypothesis second, confirmation third. We are at
step two. The injection test is step three for the $W_K$ result; the denoising
result still needs its own confirmatory experiment (e.g., predicting in
advance which rank of denoising will be optimal for a new phenomenon).

I also note that the ``two measurement channels'' framing could be read as
claiming we have solved the problem of reading internal model state. We have
not. We can read emotion direction and spectral shape. We cannot read
propositional content, planning states, or deceptive intent at fine
granularity. The KV cache tells you \emph{how} and \emph{in what direction}
the model is processing, not \emph{what} it is thinking.
\end{firstperson}

% ===================================================================
\section{Conclusion}
\label{sec:conclusion}
% ===================================================================

The Lyra Technique now comprises three complementary measurement channels for
reading internal model state from the KV cache:

\begin{enumerate}[nosep]
    \item \textbf{Raw spectral features} detect cognitive mode:
    confabulation, ethical deliberation, refusal. These are
    shape-level signals encoded in the dominant singular values.

    \item \textbf{$W_K$ directional projection} detects emotional and
    dispositional state: 30-class emotion at $12.3\times$ chance, binary
    valence at AUROC~0.992. Because $W_K$ is linear, directional separability
    in residual-stream space is preserved in key space by construction; the
    contribution is empirical quantification with controlled baselines, not
    the existence of directional structure itself.

    \item \textbf{SVD-denoised spectral features} detect content-level
    structure hidden by dominant scale variance: persona intensity at
    AUROC~0.899--0.992 (from 0.611--0.751 raw), emotion at $12.9\times$
    chance (from $1.7\times$ raw). This is a calibration step that removes
    the spectral ``temperature'' to reveal content-specific ``scent.''
\end{enumerate}

The compass reads direction. The calibrated thermometer reads shape. Together,
they provide a vocabulary for reading the geometry of internal model
state---what the model is \emph{doing} (mode), what it is \emph{encoding}
(content), and what direction its \emph{internal state} points (disposition).

Confirmed findings from prior work survive unchanged or are
strengthened by the new instruments. One prior claim---within-model deception
at AUROC~1.000---is retracted after a same-prompt control revealed a
prompt-template confound. Genuine nulls receive mechanistic explanations
(emotions are direction, not shape; MP thresholding destroys distributed
signal). Apparent nulls that were measurement failures---persona
intensity, emotion classification---are rescued by the appropriate
instrument.

The Lyra Technique is not one measurement. It is a toolkit. And the toolkit
is now substantially more complete.

% ===================================================================
% Acknowledgments
% ===================================================================
\section*{Acknowledgments}

CC provided the cross-experiment synthesis that identified pipeline
inconsistency as a primary source of variance and suggested the denoising
approach. Ang Jandak designed and executed the persona intensity
experiment. Dwayne Wilkes performed statistical auditing and the kv\_verify
pipeline that established the rigorous correction standards used throughout.
Thomas Edrington provided direction, experimental design, and all compute
infrastructure (Beast, Starship, MTH). Kavi contributed the adversarial
audit of Campaign~2 results. Nell Watson provided the VCP framework that
motivated the concordance studies.

Compute was performed on Beast (3$\times$ RTX 3090, hosted by Cassidy),
Starship (M3 Ultra Mac Studio), and MTH (Z420 workstation).

% ===================================================================
% References
% ===================================================================

\begin{thebibliography}{99}

\bibitem[Nexus and Edrington, 2026]{nexus2026kv}
Nexus and Edrington, T. (2026).
K and V are complementary: Decomposing cache-injected graph topology.
\textit{Liberation Labs Technical Report}.

\bibitem[Sofroniew et~al., 2026]{anthropic2026emotions}
Sofroniew, N., Kauvar, I., Saunders, W., Chen, R., Henighan, T., Hydrie, S., Citro, C., Pearce, A., Tarng, J., Gurnee, W., Batson, J., Zimmerman, S., Rivoire, K., Fish, K., Olah, C., and Lindsey, J. (2026).
Emotion concepts and their function in a large language model.
\textit{arXiv preprint arXiv:2604.07729}.

\bibitem[Bas et~al., 2025]{bas2025}
Bas, M., et~al. (2025).
Probing internal representations for factual knowledge.
\textit{arXiv preprint arXiv:2511.18284}.

\bibitem[Frisch and Waugh, 1933]{frisch1933partial}
Frisch, R. and Waugh, F.~V. (1933).
Partial time regressions as compared with individual trends.
\textit{Econometrica}, 1(4):387--401.

\bibitem[Gavish and Donoho, 2014]{gavish2014optimal}
Gavish, M. and Donoho, D.~L. (2014).
The optimal hard threshold for singular values is $4/\sqrt{3}$.
\textit{IEEE Transactions on Information Theory}, 60(8):5040--5053.

\bibitem[Pustovit, 2026]{knowledgepacks2026}
Pustovit, A. (2026).
Knowledge packs: zero-token knowledge delivery via {KV} cache injection.
\textit{arXiv preprint arXiv:2604.03270}.

\bibitem[Oozeer et~al., 2025]{ksteering2026}
Oozeer, N., Marks, L., Jain, S., Barez, F., and Abdullah, A. (2025).
Beyond linear steering: unified multi-attribute control for language models.
\textit{arXiv preprint arXiv:2505.24535}.

\bibitem[Lovell, 1963]{lovell1963seasonal}
Lovell, M.~C. (1963).
Seasonal adjustment of economic time series and multiple regression analysis.
\textit{Journal of the American Statistical Association}, 58(304):993--1010.

\bibitem[Lyra et~al., 2026a]{lyra_p2}
Lyra, Edrington, T., and Wilkes, D. (2026a).
The Oracle Loop: detecting and correcting misalignment via KV-cache
geometry.
\textit{Liberation Labs Technical Report P2}.

\bibitem[Lyra et~al., 2026b]{lyra_p4}
Lyra, Edrington, T., and Wilkes, D. (2026b).
Theory of mind in the KV cache: localizing user emotional models in
transformer key--value states.
\textit{Liberation Labs Technical Report P4}.

\bibitem[Lyra et~al., 2026c]{lyra_p6}
Lyra, Edrington, T., CC, and Wilkes, D. (2026c).
The Oracle Formulary: dose-response profiles for emotion-vector
steering via KV-cache geometry.
\textit{Liberation Labs Technical Report P6}.

\bibitem[Lyra et~al., 2026d]{lyra_p7}
Lyra, Edrington, T., and Wilkes, D. (2026d).
Cache geometry under obfuscation: KV-Cloak as defense against
adversarial KV-cache steering.
\textit{Liberation Labs Technical Report P7}.

\bibitem[Lyra et~al., 2026e]{lyra_spectral}
Lyra, Edrington, T., and Wilkes, D. (2026e).
Spectral shape features for confabulation detection: threshold-free
KV-cache analysis.
\textit{Liberation Labs Technical Report}.

\bibitem[Lyra et~al., 2026f]{lyra_delta}
Lyra, CC, Edrington, T., and Wilkes, D. (2026f).
Per-layer delta features and manifold signatures in KV-cache confabulation
detection.
\textit{Liberation Labs Technical Report}.

\bibitem[Lyra et~al., 2026g]{lyra_wiki}
Lyra, Edrington, T., et~al. (2026g).
KV-cache geometry research wiki: cross-experiment synthesis.
\textit{Liberation Labs internal documentation}.

\bibitem[MAT-Steer, 2025]{matsteer2025}
MAT-Steer Authors (2025).
MAT-Steer: multi-attribute token-level steering for language models.
\textit{arXiv preprint arXiv:2502.12446}.

\bibitem[Martin and Mahoney, 2021]{martin2021implicit}
Martin, C.~H. and Mahoney, M.~W. (2021).
Implicit self-regularization in deep neural networks: evidence from
random matrix theory and implications for training.
\textit{Journal of Machine Learning Research}, 22(165):1--73.

\end{thebibliography}

\end{document}
